\section{The tuned background flow}\label{sec:flow}
The linear estimates leave one expanding mean coordinate and one neutral
covariance coordinate.  We now choose the initial temperature to keep the
first coordinate small and use the averaged coefficient from
Proposition~\ref{coef:averages} to control the second.  The construction uses
estimates on finite intervals of scales; it does not assume that an infinite
bounded trajectory already exists.

Write $m_j,K_j,W_j,Y_j$ for the first four disorder cumulants of the
background list at scale $j$.  All tensor norms in this section are the
indexed projective norms of Section~\ref{sec:blocking}.  The initial bond
functions are placed in principal slots, with a crossing bond divided
symmetrically between its incident cells.  Their centered random parts are
independent symmetric two-point variables.  Consequently, uniformly for
$t=O(D^2)$,
\begin{equation}\label{flow:initial}
 g_0=\bar g(1+O(D)),\qquad \bar g=\gamma D^2,
 \qquad K_0=O(\bar g),\quad W_0=0,\quad Y_0=O(\bar g^2),
\end{equation}
where $\gamma>0$.  Positivity follows from the nonzero energy charge of a
bond and the nonzero derivative of its coupling with respect to $D$.
The initial mean is $O(|t|+D^2)$, and $x_0=x(m_0)$ has nonzero derivative
with respect to $t$ at $(t,D)=(0,0)$.  Thus $x_0$ itself can be used as the
shooting parameter on any fixed interval of size $O(D^2)$.

Put
\begin{equation}\label{flow:comparison-sequence}
 B_0=4\pi a^2\log L,\qquad
 G_j=\frac{\bar g}{1+B_0j\bar g},\qquad s_j=G_j^{1/2}.
\end{equation}
In particular $G_{j+1}/G_j=(1+B_0G_j)^{-1}$ tends uniformly to one as
$D$ tends to zero.  A trajectory is \emph{kept through scale $N$} if
$|x_j|\le XG_j$ for $0\le j\le N$, where $X$ will be fixed below.
An estimate for its tested next step concerns the actual output at $N+1$
without requiring that output to satisfy the thermal bound.

\begin{proposition}[Finite-prefix estimates]\label{flow:prefix}
There are fixed constants $X,C$ and $\rho<1$ such that, for sufficiently
small $D>0$, every trajectory kept through scale $N$ satisfies the
background tail and moment bounds of Section~\ref{sec:norms} with $s=s_j$,
and
\begin{equation}\label{flow:rough}
 \norm{m_j}+\norm{K_j}\le CG_j,\qquad
 \norm{W_j}+\norm{Y_j}\le CG_j^2,\qquad
 \tfrac12G_j<g_j<2G_j .
\end{equation}
The tensor and moment bounds also hold for the tested next step, with a
larger fixed constant for its mean.  Moreover, for $1\le j\le N$,
\begin{align}
 K_j-g_jV_j&=O\bigl(G_j(G_j+\rho^j)\bigr),\label{flow:cov-profile}\\
 m_j-g_jm_j^*&=O\bigl(G_j(G_j+\rho^j+\rho^{N-j})\bigr),
 \label{flow:mean-profile}\\
 W_j-g_j^2W_j^*,\quad Y_j-g_j^2Y_j^*
 &=O\bigl(G_j^2(s_j+\rho^j+\rho^{N-j})\bigr).
 \label{flow:higher-profile}
\end{align}
Here $V_j,m_j^*,W_j^*,Y_j^*$ are the bounded formal profiles of
Section~\ref{sec:coefficients}.  For $0\le j\le N$, including the tested
next covariance,
\begin{equation}\label{flow:variance-recursion}
 g_{j+1}=g_j+b_jg_j^2+
 O\bigl(G_j^2(s_j+\rho^j+\rho^{N-j})\bigr).
\end{equation}
The constants are independent of $N$ and of the kept trajectory.
\end{proposition}

\begin{proof}
We first derive the estimates under the corresponding bounds on an input
prefix, and then choose the constants so that the estimates reproduce
those bounds.  This separation is useful because no sign condition is
available for an individual $b_j$.

The finite Taylor rule in Proposition~\ref{norm:taylor} determines the
cumulant recursion with remainder $O(G_j^{5/2})$.  The required algebra can
be described explicitly.  Assign a separate letter to every factor in a
Taylor monomial, even when two letters have the same slot index.  If
$I_1,\ldots,I_k$ are the groups belonging to the $k$ output factors, the
cumulant of their products is
\begin{equation}\label{flow:connected-partitions}
 \cum\left(\prod_{a\in I_1}h_a,\ldots,
                 \prod_{a\in I_k}h_a\right)
 =\sum_{\pi:\,\pi\vee\{I_1,\ldots,I_k\}=\one}
       \ \prod_{B\in\pi}\cum(h_a:a\in B).
\end{equation}
The join condition says precisely that the partition connects the output
groups.  Formula~\eqref{flow:connected-partitions} follows by expanding
moments into cumulants and applying the defining alternating sum for the
output cumulant.  The coefficient of a partition not connecting the groups
is zero.  It therefore applies also to correlated input slots.

A singleton block costs $m=O(G)$, a pair costs $K=O(G)$, and blocks of
sizes three and four cost $W,Y=O(G^2)$, where $G=G_j$.  For covariance the
nonlinear degree patterns with at most four letters are
$(1,2),(2,1),(1,3),(2,2),(3,1)$.
The three-letter patterns retain $W$ and the connecting $mK$ partitions;
the four-letter patterns retain $Y$ and the connecting $KK$ partitions.
For the third cumulant the only further pattern has four letters and
retains $Y$ and $KK$.  For the fourth cumulant there is only homogeneous
linear transport to the stated accuracy.  The mean retains its linear
term and the quadratic Taylor coefficient applied to $K$; its other
terms have size $O(G^2)$.  In particular the recursions have the form
\begin{align}
 m_{j+1}&=T_jm_j+Q_j(K_j)+O(G_j^2),\label{flow:mean-recursion}\\
 K_{j+1}&=T_j^{(2)}K_j+
       \mathcal B_j(W_j,Y_j,K_j\otimes K_j,m_j\otimes K_j)
       +O(G_j^{5/2}),\label{flow:cov-recursion}\\
 W_{j+1}&=T_j^{(3)}W_j+
       \mathcal F_j(Y_j,K_j\otimes K_j)+O(G_j^{5/2}),
       \label{flow:third-recursion}\\
 Y_{j+1}&=T_j^{(4)}Y_j+O(G_j^{5/2}).\label{flow:fourth-recursion}
\end{align}
Here $Q_j$ is the quadratic Taylor coefficient contracted against a
covariance; $\mathcal B_j$ and $\mathcal F_j$ denote the sums of exactly
the connected partitions just listed, including their Taylor
coefficients.  They are bounded linear maps on the displayed direct
sums of tensor spaces.  The indexed estimates in
Proposition~\ref{norm:taylor} justify all slot sums.

The charge of linear covariance transport is unchanged by
Proposition~\ref{lin:background}.  Equations~\eqref{flow:rough}
and~\eqref{flow:cov-recursion} therefore first give the coarse estimate
$g_{j+1}-g_j=O(G_j^2)$.  On the kernel of that charge the linear operator
contracts.  Subtracting $g_{j+1}V_{j+1}$, and using the exact transport of
$V_j$, gives~\eqref{flow:cov-profile} by convolution with this contraction.
For clarity, every such convolution uses the elementary estimate
\begin{equation}\label{flow:convolution}
 \sum_{i<j}\kappa^{j-1-i}G_i^\alpha
 \le C_\alpha G_j^\alpha,\qquad
 \kappa^jG_0^\alpha\le C_\alpha\rho^jG_j^\alpha,
\end{equation}
for the finitely many positive exponents $\alpha$ used here, where
$\kappa<\rho<1$.  To verify it, use
$G_i/G_j\le(1+B_0\bar g)^{j-i}$ and decrease $D$ so that
$\kappa(1+B_0\bar g)^\alpha<\rho$.

Subtract next the recursion for $g_jm_j^*$ from
\eqref{flow:mean-recursion}.  The stable mean component is controlled
forward by~\eqref{flow:convolution}.  Its forcing is
$O(G_j^2+G_j\rho^j)$.  The thermal component is controlled backward:
if its error is $z_j$, then
\[
 z_j=L^{-1}z_{j+1}+O(G_j^2+G_j\rho^j),\qquad |z_N|\le CG_N.
\]
Summing this equality from $N$ backwards gives
$|z_j|\le CG_j(G_j+\rho^j+\rho^{N-j})$, after enlarging $\rho<1$.
This proves~\eqref{flow:mean-profile}.  The fourth and third cumulant
recursions, in that order, have contracting homogeneous operators.
Subtract their formal-profile recursions, use
$g_{j+1}^2-g_j^2=O(G_j^3)$, and apply
\eqref{flow:convolution}.  The resulting errors are those in
\eqref{flow:higher-profile}.  An unmatched initial profile contributes
only the geometric term.  All comparisons at scale zero can likewise
be paid by that term.

The definition of $b_j$ is the charge of the nonlinear term in
\eqref{flow:cov-recursion} evaluated on the formal profiles.
Inserting~\eqref{flow:cov-profile}--\eqref{flow:higher-profile}
therefore gives~\eqref{flow:variance-recursion}.  The terminal term in
that formula is the cost of using only a finite thermal prefix.

We now show why this refined estimate reproduces the asserted covariance
range.  The coarse one-step bound first makes the tested $g_{N+1}$
positive and comparable to $G_{N+1}$, so reciprocals are legitimate.
Taking reciprocals in~\eqref{flow:variance-recursion} and summing gives,
for $1\le k\le N+1$,
\begin{equation}\label{flow:reciprocal}
 \frac1{g_k}=\frac1{g_0}-\sum_{j<k}b_j+
 O\left(1+\sum_{j<k}s_j\right).
\end{equation}
Indeed $g_j/G_j$ is comparable to one on the input prefix, the two
geometric errors have bounded sums, and the reciprocal Taylor remainder
is $O(G_j)$, bounded by $O(s_j)$.
Proposition~\ref{coef:averages} says that
$R_k=\sum_{j<k}(b_j+B_0)=o(k)$.  Thus for each $\eta>0$ there is
$C_\eta<\infty$ such that $|R_k|\le\eta k+C_\eta$ for every $k$.
Since $s_j\le\sqrt{\bar g}$,
\begin{equation}\label{flow:relative-reciprocal}
 \sup_{1\le k\le N+1}
 \frac{\left|g_k^{-1}-(\bar g^{-1}+B_0k)\right|}
      {\bar g^{-1}+B_0k}
 \le C D+C\sqrt{\bar g}+C\eta+C_\eta\bar g.
\end{equation}
First choosing $\eta$ small and then $D$ small makes this strictly
smaller than the amount allowed by $G_k/2<g_k<2G_k$.
It also proves that
\begin{equation}\label{flow:uniform-ratio}
 \sup_{k\le N+1}|g_k/G_k-1|\longrightarrow0\quad(D\downarrow0),
\end{equation}
uniformly over all kept prefixes.  Only the averaged coefficient was
used; individual coefficients may have either sign.

It remains to fix the constants in the bounds used above.  First choose
the covariance-kernel allowance larger than its initial and line
contributions.  Contraction reproduces this allowance, with its nonlinear
$O(G^2)$ loss left for the final smallness choice.  Choose $X$ so large
that, under this covariance bound,
\begin{equation}\label{flow:outward}
 x_{j+1}=Lx_j+f_j,\qquad |f_j|\le C G_j,
 \qquad (L-1)X>C+1.
\end{equation}
Choose the remaining mean allowance from this tube, the initial mean,
and stable contraction.  Next choose the fourth-cumulant allowance and
then the third-cumulant allowance, using
\eqref{flow:fourth-recursion} and~\eqref{flow:third-recursion}.
These choices furnish a fixed bound $F_*$ on the actual next principal
raw second tensor.  Indeed, at each principal slot $i$,
\[
 \norm{\E[h'_i\otimes h'_i]}
 =\norm{K'_{i,i}+m'_i\otimes m'_i}
 \le C G_j+C G_j^2\le e^{-2v}F_*^2G_j,
\]
after enlarging $F_*$ by the fixed factor $e^v$ if necessary.
This is precisely the tensor input required by the moment reproduction
estimate; no output probability norm has been assumed to obtain it.

For the moment bounds in Proposition~\ref{norm:reproduction}, choose
$B_2$ larger than the fixed multiple of $F_*$ and the initial second
moment bound, and $B$ larger than the initial high-moment bound and
$C_L^*B_2$, with strict room in both inequalities.  The dictionary error
was chosen to have $\eps_* C_L^*$ sufficiently small, so these
choices can also satisfy $\eps_*B\ll B_2$.  Finally decrease $D$ to
pay all nonlinear remainders, the variation of $G_j$, and all fixed
constants in the profile comparisons.  The tail bound is reproduced by
the exact map.  This order closes the input-prefix induction and its
tested next step.  In particular, the covariance improvement uses only
the previously kept prefix, not an as yet unproved next thermal bound.
\end{proof}

\begin{proposition}[Existence of a tuned trajectory]\label{flow:tuned}
For every sufficiently small $D>0$ there is $t_*(D)=O(D^2)$ whose
background trajectory satisfies $|x_j|\le XG_j$ at every scale.
All the bounds of Proposition~\ref{flow:prefix} hold along this
trajectory.  In the profile estimates its terminal geometric term may
be omitted.
\end{proposition}

\begin{proof}
At every fixed scale the expected list coordinates are continuous in
$x_0$.  Indeed the map has finite random footprints; each fresh threshold
has a continuous distribution, so equality with that threshold has
probability zero.  Almost-sure continuity away from those equalities,
together with the finite-step moment bounds and summable coefficient
bounds, gives dominated convergence.  This applies uniformly on each
compact parameter interval used below.

Choose the initial interval $I_0$ so that its endpoints have
$x_0=-XG_0$ and $x_0=XG_0$.  Suppose $I_j$ is a closed interval all of
whose parameters are kept through $j$, and whose endpoints have
$x_j=-XG_j$ and $x_j=XG_j$.  By~\eqref{flow:outward}, their next
coordinates lie strictly below $-XG_{j+1}$ and above $XG_{j+1}$,
respectively.  Continuity supplies a subinterval $I_{j+1}\subset I_j$
on which $|x_{j+1}|\le XG_{j+1}$ and whose endpoints attain both bounds.
For an explicit choice, take the first hitting of the upper bound and
the last preceding hitting of the lower bound.  Between those points
neither bound can be crossed.  The nested nonempty compact intervals
have a common parameter, giving the asserted trajectory.  Finally,
apply the finite-prefix profile bounds to arbitrarily long prefixes of
this trajectory, with $j$ fixed, and let $N$ tend to infinity.
\end{proof}

\begin{corollary}[Decay of the disorder coordinate]\label{flow:variance}
Along every tuned trajectory obtained above,
\begin{equation}\label{flow:variance-asymptotic}
 g_j\sim\frac1{B_0j}\quad(j\to\infty),\qquad
 \sum_{j\ge0}|g_{j+1}-g_j|\le C\bar g.
\end{equation}
The uniform comparison~\eqref{flow:uniform-ratio} also holds over all
$j\ge0$ as $D\downarrow0$.
\end{corollary}

\begin{proof}
In~\eqref{flow:reciprocal}, $R_j=o(j)$ and
$\sum_{i<j}s_i=o(j)$ for each fixed $D>0$, because $s_i\to0$.
Thus $g_j^{-1}=B_0j+o(j)$.  The coarse increment estimate and
$\sum_jG_j^2\le C\bar g$ give the total-variation bound.
\end{proof}
